\SourcePage{14}{19}
\section{Quantization of the Free Electromagnetic Field}

In order to treat the electromagnetic field as a quantum object, it is convenient to start from a classical description in which the field is characterized by a discrete, albeit infinite, set of variables; such a description will allow one to apply the standard formalism of quantum mechanics directly. A description by means of potentials specified at every point of space is, in essence, a description in terms of a continuous set of variables.

Let $\mathbf{A}(\mathbf{r},t)$ be the vector potential of the free electromagnetic field, satisfying the ``transversality condition''
\begin{equation}
\Div{\mathbf{A}} = 0.
\tag{2.1}
\end{equation}
The scalar potential then vanishes, $\Phi = 0$, and the fields $\mathbf{E}$ and $\mathbf{H}$ are
\begin{equation}
\mathbf{E} = -\dot{\mathbf{A}}, \qquad \mathbf{H} = \Curl{\mathbf{A}}.
\tag{2.2}
\end{equation}
Maxwell's equations reduce to the wave equation for~$\mathbf{A}$:
\begin{equation}
\Delta\mathbf{A} - \frac{\partial^2 \mathbf{A}}{\partial t^2} = 0.
\tag{2.3}
\end{equation}

As is well known (see~Vol.~II, \S~52), in classical electrodynamics the passage to a description in terms of a discrete set of variables is effected by considering the field inside a large but finite volume~$V$.\footnote{To avoid burdening formulae with superfluous factors, we set $V = 1$.} Let us recall, omitting the details of the calculation, how this is done.

A field in a finite volume can be expanded in travelling plane waves, so that its potential takes the form of a series
\begin{equation}
\mathbf{A} = \sum_{\mathbf{k}} (\mathbf{a}_{\mathbf{k}} e^{i\mathbf{kr}} + \mathbf{a}_{\mathbf{k}}^* e^{-i\mathbf{kr}}),
\tag{2.4}
\end{equation}
where the coefficients $\mathbf{a}_{\mathbf{k}}$ depend on time according to
\begin{equation}
\mathbf{a}_{\mathbf{k}} \sim e^{-i\omega t}, \quad \omega = |\mathbf{k}|.
\tag{2.5}
\end{equation}
By virtue of condition~(2.1) the complex vectors $\mathbf{a}_{\mathbf{k}}$ are orthogonal to the corresponding wave vectors: $\mathbf{a}_{\mathbf{k}} \mathbf{k} = 0$.
\SourcePage{15}{20}

The summation in~(2.4) runs over an infinite discrete set of values of the wave vector (its three components $k_x$, $k_y$, $k_z$). The transition to an integration over a continuous distribution is made with the aid of the expression
\[
\frac{d^3k}{(2\pi)^3}
\]
for the number of allowed values of~$\mathbf{k}$ falling within a volume element $d^3k = dk_x\,dk_y\,dk_z$ of $\mathbf{k}$-space.

Once the vectors~$\mathbf{a}_{\mathbf{k}}$ are given, the field inside the volume in question is completely determined. One may therefore view these quantities as a discrete set of classical ``field variables''. In order to find the passage to a quantum theory, however, a further transformation of these quantities is needed, one that brings the field equations into a form analogous to the canonical (Hamiltonian) equations of classical mechanics. The canonical variables of the field are introduced by means of
\begin{equation}
\mathbf{Q}_{\mathbf{k}} = \frac{1}{\sqrt{4\pi}}(\mathbf{a}_{\mathbf{k}} + \mathbf{a}_{\mathbf{k}}^*), \qquad
\mathbf{P}_{\mathbf{k}} = -\frac{i\omega}{\sqrt{4\pi}}(\mathbf{a}_{\mathbf{k}} - \mathbf{a}_{\mathbf{k}}^*) = \dot{\mathbf{Q}}_{\mathbf{k}}
\tag{2.6}
\end{equation}
(they are evidently real). The vector potential is expressed through the canonical variables as
\begin{equation}
\mathbf{A} = \sqrt{4\pi} \sum_{\mathbf{k}} \biggl(\mathbf{Q}_{\mathbf{k}} \cos\mathbf{kr} - \frac{1}{\omega}\mathbf{P}_{\mathbf{k}} \sin\mathbf{kr}\biggr).
\tag{2.7}
\end{equation}

To find the Hamiltonian~$H$ one must evaluate the total field energy
\[
\frac{1}{8\pi}\int (\mathbf{E}^2 + \mathbf{H}^2)\,d^3x,
\]
expressed in terms of the quantities~$\mathbf{Q}_{\mathbf{k}}$, $\mathbf{P}_{\mathbf{k}}$. Substituting the expansion~(2.7) for~$\mathbf{A}$, computing $\mathbf{E}$ and $\mathbf{H}$ by means of~(2.2), and integrating, one obtains
\[
H = \frac{1}{2}\sum_{\mathbf{k}} (\mathbf{P}_{\mathbf{k}}^2 + \omega^2 \mathbf{Q}_{\mathbf{k}}^2).
\]

Each of the vectors $\mathbf{P}_{\mathbf{k}}$ and $\mathbf{Q}_{\mathbf{k}}$ is perpendicular to the wave vector~$\mathbf{k}$ and therefore has two independent components. The direction of these vectors fixes the polarization direction of the corresponding wave. Denoting the two components of $\mathbf{Q}_{\mathbf{k}}$, $\mathbf{P}_{\mathbf{k}}$ (in the plane perpendicular to~$\mathbf{k}$) by $Q_{\mathbf{k}\alpha}$, $P_{\mathbf{k}\alpha}$ ($\alpha = 1, 2$), we rewrite the Hamiltonian as
\begin{equation}
H = \sum_{\mathbf{k}\alpha} \frac{1}{2} (P_{\mathbf{k}\alpha}^2 + \omega^2 Q_{\mathbf{k}\alpha}^2).
\tag{2.8}
\end{equation}
\SourcePage{16}{21}

The Hamiltonian has thus broken up into a sum of independent terms, each involving just one pair $Q_{\mathbf{k}\alpha}$, $P_{\mathbf{k}\alpha}$. Every term answers to a travelling wave of a given wave vector and polarization and reproduces the structure of a single harmonic-oscillator Hamiltonian. One accordingly speaks of this decomposition as the expansion of the field in \emph{oscillators}.

We now turn to the quantization of the free electromagnetic field. The method of classical description just set out makes the route to quantum theory obvious. The canonical variables~--- the generalized coordinates $Q_{\mathbf{k}\alpha}$ and generalized momenta $P_{\mathbf{k}\alpha}$~--- are now to be treated as operators obeying the commutation rule
\begin{equation}
\hat{P}_{\mathbf{k}\alpha}\hat{Q}_{\mathbf{k}\alpha} - \hat{Q}_{\mathbf{k}\alpha}\hat{P}_{\mathbf{k}\alpha} = -i
\tag{2.9}
\end{equation}
(operators with different $\mathbf{k}\alpha$ all commute with one another). Together with them the potential~$\mathbf{A}$ and, by~(2.2), the field strengths $\mathbf{E}$ and~$\mathbf{H}$ also become (Hermitian) operators.

A consistent determination of the field Hamiltonian requires evaluation of the integral
\begin{equation}
\hat{H} = \frac{1}{8\pi}\int (\hat{\mathbf{E}}^2 + \hat{\mathbf{H}}^2)\,d^3x,
\tag{2.10}
\end{equation}
in which $\hat{\mathbf{E}}$ and $\hat{\mathbf{H}}$ are expressed through the operators $\hat{P}_{\mathbf{k}\alpha}$, $\hat{Q}_{\mathbf{k}\alpha}$. In practice, however, the non-commutativity of the latter does not make itself felt here, since the products $Q_{\mathbf{k}\alpha} P_{\mathbf{k}\alpha}$ enter with the factor $\cos\mathbf{kr}\cdot\sin\mathbf{kr}$, which vanishes upon integration over the full volume. Consequently the Hamiltonian takes the form
\begin{equation}
\hat{H} = \sum_{\mathbf{k}\alpha} \frac{1}{2}(\hat{P}_{\mathbf{k}\alpha}^2 + \omega^2 \hat{Q}_{\mathbf{k}\alpha}^2),
\tag{2.11}
\end{equation}
exactly reproducing the classical Hamiltonian, as one would naturally expect.

Finding the eigenvalues of this Hamiltonian requires no special computation, since the problem reduces to the familiar one of the energy levels of linear oscillators (see~Vol.~III, \S~23). We can therefore write down at once the energy levels of the field:
\begin{equation}
E = \sum_{\mathbf{k}\alpha}\biggl(N_{\mathbf{k}\alpha} + \frac{1}{2}\biggr)\omega,
\tag{2.12}
\end{equation}
where $N_{\mathbf{k}\alpha}$ are non-negative integers.

Discussion of this formula will be taken up in the next section; here we write out the matrix elements of the quantities $Q_{\mathbf{k}\alpha}$,
\SourcePage{17}{22}
which can be obtained directly from the known formulae for the matrix elements of the oscillator coordinate (see~Vol.~III, \S~23). The non-vanishing matrix elements are
\begin{equation}
\langle N_{\mathbf{k}\alpha} | Q_{\mathbf{k}\alpha} | N_{\mathbf{k}\alpha} - 1\rangle = \sqrt{\frac{N_{\mathbf{k}\alpha}}{2\omega}}.
\tag{2.13}
\end{equation}

The matrix elements of $P_{\mathbf{k}\alpha} = \dot{Q}_{\mathbf{k}\alpha}$ differ from those of $Q_{\mathbf{k}\alpha}$ only by a factor $\pm i\omega$.

For subsequent calculations it will, however, be more convenient to work not with $Q_{\mathbf{k}\alpha}$, $P_{\mathbf{k}\alpha}$ themselves but with their linear combinations $\omega Q_{\mathbf{k}\alpha} \pm iP_{\mathbf{k}\alpha}$, which have matrix elements only for the transitions $N_{\mathbf{k}\alpha} \to N_{\mathbf{k}\alpha} \pm 1$. Accordingly we introduce the operators
\begin{equation}
\hat{c}_{\mathbf{k}\alpha} = \frac{1}{\sqrt{2\omega}}(\omega\hat{Q}_{\mathbf{k}\alpha} + i\hat{P}_{\mathbf{k}\alpha}), \qquad
\hat{c}_{\mathbf{k}\alpha}^+ = \frac{1}{\sqrt{2\omega}}(\omega\hat{Q}_{\mathbf{k}\alpha} - i\hat{P}_{\mathbf{k}\alpha})
\tag{2.14}
\end{equation}
(the classical quantities $c_{\mathbf{k}\alpha}$, $c_{\mathbf{k}\alpha}^*$ coincide, up to a factor $\sqrt{2\pi/\omega}$, with the coefficients $a_{\mathbf{k}\alpha}$, $a_{\mathbf{k}\alpha}^*$ in the expansion~(2.4)).

The matrix elements of these operators are
\begin{equation}
\langle N_{\mathbf{k}\alpha} - 1 | \hat{c}_{\mathbf{k}\alpha} | N_{\mathbf{k}\alpha}\rangle = \langle N_{\mathbf{k}\alpha} | \hat{c}_{\mathbf{k}\alpha}^+ | N_{\mathbf{k}\alpha} - 1\rangle = \sqrt{N_{\mathbf{k}\alpha}}.
\tag{2.15}
\end{equation}

The commutation rule between $\hat{c}_{\mathbf{k}\alpha}$ and $\hat{c}_{\mathbf{k}\alpha}^+$ follows from the definition~(2.14) and the rule~(2.9):
\begin{equation}
\hat{c}_{\mathbf{k}\alpha}\hat{c}_{\mathbf{k}\alpha}^+ - \hat{c}_{\mathbf{k}\alpha}^+\hat{c}_{\mathbf{k}\alpha} = 1.
\tag{2.16}
\end{equation}

Returning to the vector potential, we adopt an expansion of the same type as~(2.4), but with the coefficients now promoted to operators:
\begin{equation}
\hat{\mathbf{A}} = \sum_{\mathbf{k}\alpha} (\hat{c}_{\mathbf{k}\alpha}\mathbf{A}_{\mathbf{k}\alpha} + \hat{c}_{\mathbf{k}\alpha}^+\mathbf{A}_{\mathbf{k}\alpha}^*),
\tag{2.17}
\end{equation}
where
\begin{equation}
\mathbf{A}_{\mathbf{k}\alpha} = \sqrt{4\pi}\,\frac{\mathbf{e}^{(\alpha)}}{\sqrt{2\omega}}\, e^{i\mathbf{kr}}.
\tag{2.18}
\end{equation}

Here $\mathbf{e}^{(\alpha)}$ denotes a unit vector specifying the polarization direction of the oscillator; the vectors $\mathbf{e}^{(\alpha)}$ are perpendicular to the wave vector~$\mathbf{k}$, and for each~$\mathbf{k}$ there are two independent polarizations.

Analogously, for the operators $\hat{\mathbf{E}}$ and $\hat{\mathbf{H}}$ we write
\begin{equation}
\hat{\mathbf{E}} = \sum_{\mathbf{k}\alpha} (\hat{c}_{\mathbf{k}\alpha}\mathbf{E}_{\mathbf{k}\alpha} + \hat{c}_{\mathbf{k}\alpha}^+\mathbf{E}_{\mathbf{k}\alpha}^*), \quad
\hat{\mathbf{H}} = \sum_{\mathbf{k}\alpha} (\hat{c}_{\mathbf{k}\alpha}\mathbf{H}_{\mathbf{k}\alpha} + \hat{c}_{\mathbf{k}\alpha}^+\mathbf{H}_{\mathbf{k}\alpha}^*),
\tag{2.19}
\end{equation}
with
\begin{equation}
\mathbf{E}_{\mathbf{k}\alpha} = i\omega\mathbf{A}_{\mathbf{k}\alpha}, \qquad \mathbf{H}_{\mathbf{k}\alpha} = \mathbf{n}\times\mathbf{E}_{\mathbf{k}\alpha} \qquad (\mathbf{n} = \mathbf{k}/\omega).
\tag{2.20}
\end{equation}
\SourcePage{18}{23}

The vectors $\mathbf{A}_{\mathbf{k}\alpha}$ are mutually orthogonal in the sense that
\begin{equation}
\int \mathbf{A}_{\mathbf{k}\alpha}\mathbf{A}_{\mathbf{k}'\alpha'}^*\,d^3x = \frac{2\pi}{\omega}\,\delta_{\alpha\alpha'}\,\delta_{\mathbf{kk}'}.
\tag{2.21}
\end{equation}
To see this, note that when $\mathbf{A}_{\mathbf{k}\alpha}$ and $\mathbf{A}_{\mathbf{k}'\alpha'}^*$ carry different wave vectors their product involves $e^{i(\mathbf{k}-\mathbf{k}')\mathbf{r}}$, which integrates to zero over the volume; when only the polarizations differ, the orthogonality $\mathbf{e}^{(\alpha)}\mathbf{e}^{(\alpha')*} = 0$ of the two independent polarization directions provides the required vanishing. Analogous relations hold for the vectors $\mathbf{E}_{\mathbf{k}\alpha}$, $\mathbf{H}_{\mathbf{k}\alpha}$. Their normalization is conveniently written as
\begin{equation}
\frac{1}{4\pi}\int (\mathbf{E}_{\mathbf{k}\alpha}\mathbf{E}_{\mathbf{k}'\alpha'}^* + \mathbf{H}_{\mathbf{k}\alpha}\mathbf{H}_{\mathbf{k}'\alpha'}^*)\,d^3x = \omega\,\delta_{\mathbf{kk}'}\,\delta_{\alpha\alpha'}.
\tag{2.22}
\end{equation}

Substituting the operators~(2.19) into~(2.10) and integrating with the help of~(2.22), we get the field Hamiltonian expressed through the operators $\hat{c}$, $\hat{c}^+$:
\begin{equation}
\hat{H} = \sum_{\mathbf{k}\alpha} \frac{1}{2}\omega(\hat{c}_{\mathbf{k}\alpha}\hat{c}_{\mathbf{k}\alpha}^+ + \hat{c}_{\mathbf{k}\alpha}^+\hat{c}_{\mathbf{k}\alpha}).
\tag{2.23}
\end{equation}
In the representation under discussion (with the matrix elements of $\hat{c}$, $\hat{c}^+$ given by~(2.15)) this operator is diagonal, and its eigenvalues coincide, naturally, with~(2.12).

In classical theory the field momentum is defined by the integral
\[
\mathbf{P} = \frac{1}{4\pi}\int \mathbf{E}\times\mathbf{H}\,d^3x.
\]
On passing to the quantum theory by replacing $\mathbf{E}$ and $\mathbf{H}$ with the operators~(2.19), one readily finds
\begin{equation}
\hat{\mathbf{P}} = \sum_{\mathbf{k}\alpha} \frac{1}{2}(\hat{P}_{\mathbf{k}\alpha}^2 + \omega^2\hat{Q}_{\mathbf{k}\alpha}^2)\mathbf{n}
\tag{2.24}
\end{equation}
--- which mirrors the familiar classical link between energy and momentum of plane waves. This operator has eigenvalues
\begin{equation}
\mathbf{P} = \sum_{\mathbf{k}\alpha} \mathbf{k}\biggl(N_{\mathbf{k}\alpha} + \frac{1}{2}\biggr).
\tag{2.25}
\end{equation}

The operator representation realized by the matrix elements~(2.15) is called the ``occupation-number representation'': it amounts to characterizing the state of the system (the field) through the quantum numbers~$N_{\mathbf{k}\alpha}$ (\emph{occupation numbers}). Within this representation the field operators~(2.19) (and hence the Hamiltonian~(2.11)) act on the system's wave function, written as a function of the
\SourcePage{19}{24}
numbers~$N_{\mathbf{k}\alpha}$; we denote it $\Phi(N_{\mathbf{k}\alpha},t)$. The field operators~(2.19) do not depend on time explicitly. This corresponds to the Schr\"odinger representation of operators, familiar from non-relativistic quantum mechanics. What does depend on time is the state of the system, $\Phi(N_{\mathbf{k}\alpha},t)$, its time dependence being governed by the Schr\"odinger equation
\[
i\frac{\partial\Phi}{\partial t} = \hat{H}\Phi.
\]

This description of the field is relativistically invariant in essence, since it rests on the invariant Maxwell equations. Yet the invariance is not exhibited explicitly~--- above all because the spatial coordinates and the time enter the description in a highly asymmetric way.

Within a relativistic framework it is desirable to give the description a more overtly invariant appearance. For this purpose one uses the so-called Heisenberg representation, wherein the explicit dependence on time is shifted onto the operators (see~Vol.~III, \S~13). Time and spatial coordinates then appear symmetrically in the field-operator expressions, while the system state $\Phi$ depends only on the occupation numbers.

For the operator $\hat{\mathbf{A}}$ the passage to the Heisenberg representation amounts to replacing, in each term of the sum in~(2.17), (2.18), the factor $e^{i\mathbf{kr}}$ by $e^{i(\mathbf{kr}-\omega t)}$; that is, one is to understand by~$\mathbf{A}_{\mathbf{k}\alpha}$ the time-dependent functions
\begin{equation}
\mathbf{A}_{\mathbf{k}\alpha} = \sqrt{4\pi}\,\frac{\mathbf{e}^{(\alpha)}}{\sqrt{2\omega}}\, e^{-i(\omega t - \mathbf{kr})}.
\tag{2.26}
\end{equation}

This is easily verified by noting that the matrix element of a Heisenberg operator for the transition $i \to f$ must contain the factor $\exp\{-i(E_i - E_f)t\}$, where $E_i$ and $E_f$ are the energies of the initial and final states (see~Vol.~III, \S~13). For a transition in which~$N_{\mathbf{k}}$ decreases or increases by one, this factor reduces to $e^{-i\omega t}$ or $e^{i\omega t}$, respectively. The indicated replacement satisfies this requirement.

Henceforth (in treating both the electromagnetic field and the fields of particles) we shall always understand the Heisenberg representation of operators.
